% 進捗報告テンプレート（1報告＝1ページ。REP1つにつき1本）
% 1ページ目：仮説検証型（検証）／2ページ目：探索型（探索）。
% 使うときは，使わない型のブロックと \clearpage を削除し，REPごとに1ファイルにする。
% IDは自動では振らない。\reportid に REP のID（例 REP003），\questionid に問いのID（例 Q001-A）を書く。
% 書き方のルールは .claude/commands/make_progress_report.md，置き場所とIDは research/README.md。
% コンパイル：uplatex を2回 → dvipdfmx。progress_report.sty は TEXINPUTS で
% .claude/templates/progress_report// を参照する（コピーしない）。
\documentclass[twocolumn, 10pt]{ujarticle}
\usepackage{progress_report}

\begin{document}

% ################################################################
% 仮説検証型（検証したい仮説が1つあるとき）
% ################################################################
\reportid{REP000}
\questionid{Q000-A}
\verification{ここに検証のタイトルを書く}
\reportdate{20261006}
\summary{%
ここに報告全体のサマリーを2〜3文の平文で書く。
1文目に検証した仮説，2文目に検証方法，3文目に結果と判断を書く。
実験が未実施であれば，3文目は「実験は未実施である」と書く。
}
\beginverification

% ============================================================
\section{目的}\label{sec:purpose}
% ============================================================

この検証で何を明らかにしたいかを，簡潔に1〜2文で書く。手法の性能ではなく，
どの説明を確かめるのかだけを書き，観測した事実はここに書かない。前の報告がある場合は，その今後の展望のどれを
受けた目的かが分かるように書き出す。

% ============================================================
\section{実験前に考えた仮説}\label{sec:prehyp}
% ============================================================

\subsection{観測した事実}\label{sec:fact}

仮説の根拠になった事実だけを，簡潔に書く。条件・図表・数値をそのまま
記載し，解釈はここに書かない。前の報告から引き継いだ事実には，その報告のID（例 REP002）を付ける。

% 図の例
% \begin{figure}[H]
% \centering
% \includegraphics[width=\linewidth]{fig_xxx.png}
% \caption{図の説明}
% \label{fig:xxx}
% \end{figure}

% 表の例
% \begin{table}[H]
% \centering
% \caption{表の説明}
% \label{tab:xxx}
% \footnotesize
% \begin{tabularx}{\linewidth}{@{}Xrr@{}}
% \toprule
% 条件 & 指標A & 指標B \\
% \midrule
% 条件1 & 0.00 & 0.00 \\
% 条件2 & 0.00 & 0.00 \\
% \bottomrule
% \end{tabularx}
% \end{table}

\subsection{事実から考えた仮説}\label{sec:hyp}

最初に，\S\ref{sec:fact}の事実について何を知りたいのかを平文で書く。

\hypothesis{今回検証する仮説を1つ，1文で書く。}

枠の下に，なぜそう考えたかを，事実から仮説に至る筋道が追えるように丁寧に
書く。見出しやラベルは付けず，論文の本文のように文章で続ける。事実と，
そこから私が推測したことを，文を分けて書く。式は仮説の根拠であり，
検証の結果ではないことを明示する。同じ事実を説明できる別の原因が実際に
あるときは，その説明と，それが正しい場合に結果がどうなるかも書く。

\subsection{仮説に基づく予測}\label{sec:pred}

何を変えると，何がどう変わるはずかを，実験の前に書く。最終的な性能だけでなく，
仮説が正しいなら同時に変わるはずの量も挙げる。

% ============================================================
\section{実験}\label{sec:exp}
% ============================================================

\subsection{検証方法}\label{sec:method}

固定する条件と変える条件を，簡潔に表で書く。

\subsection{判断指標}\label{sec:criteria}

測定する量と，それをどう読んで判断するかを，実験の前に決めて丁寧に書く。
「支持された」「支持されなかった」「まだ判断できない」のそれぞれが，どの
結果に当たるかを具体的な数値や向きで書く。仮説を支持しない結果は，
特に明記する。

% ============================================================
\section{結果}\label{sec:result}
% ============================================================

実験で観測した事実だけを，簡潔に書く。\S\ref{sec:criteria}で決めた
測定量を，同じ順に並べる。原因の説明や評価の言葉はここに書かない。

% ============================================================
\section{考察}\label{sec:disc}
% ============================================================

最初に，\S\ref{sec:criteria}の判断指標に照らして「支持された」「支持されなかった」
「まだ判断できない」のいずれかを書く。続けて，丁寧に書く。
結果から言えることには「〜であった」「〜が示された」，私の解釈には
「〜と考えられる」「〜の可能性がある」を使い，文を分ける。
どこまで言えるか，同じ結果を説明できる別の説明は何か，何がまだ
分からないかを書く。支持されなかった場合に，追加の仮説で説明して
終わらせない。

% ============================================================
\section{今後の展望}\label{sec:future}
% ============================================================

次の報告の目的に繋がるように書く。ラベルや箇条書きは使わず，3つの段落で，
この順に書く。まず，今回の結果から確かに言えることを，「〜であった」の文末で
書く。次に，まだ確かめていないが今考えていることを，「〜と考えられる」の文末で
書く。最後に，次に知りたいことを，「〜を知りたい」の文で1つ書く。

\clearpage

% ################################################################
% 探索型（仮説がないとき。動くか確認したい／新しすぎて仮説が立たない／
%         仮説を立てるために現象を見たい）
% ################################################################
\reportid{REP000}
\questionid{Q000-A}
\exploration{ここに確認したい現象をタイトルに書く}
\reportdate{20261006}
\summary{%
ここに報告全体のサマリーを2〜3文の平文で書く。
1文目に確認したかった現象，2文目に確認の方法，3文目に観測した結果と，そこから生まれた問いを書く。
実験が未実施であれば，3文目は「実験は未実施である」と書く。
}
\beginverification

% ============================================================
\section{目的}\label{sec:e-purpose}
% ============================================================

何を知りたいかを，1〜2文で書く。今の状況や確認項目の一覧は，ここに書かない（REPのmdに書く）。

% ============================================================
\section{実験}\label{sec:e-exp}
% ============================================================

\subsection{対象のアルゴリズム}\label{sec:e-algo}

手法に共通すること（何を決めるのか，手法どうしで何が違うのか）を，1〜2文で書く。手法ごとの中身は表にする。

\begin{table}[H]
\centering
\caption{対象の手法（数値は既定値）}
\label{tab:e-algo}
\footnotesize
\begin{tabularx}{\linewidth}{@{}l>{\raggedright\arraybackslash}p{0.2\linewidth}>{\raggedright\arraybackslash}X@{}}
\toprule
手法 & 見るもの & 決め方 \\
\midrule
手法A & 入力を短く & 決め方を1〜2行で \\
\addlinespace
手法B & 入力を短く & 決め方を1〜2行で \\
\bottomrule
\end{tabularx}
\end{table}

\subsection{実験環境}\label{sec:e-env}

評価に使う量の定義を，1〜2文で書く。固定した条件と変えた条件は表にする。セルは1行1項目で短く書く。

\begin{table}[H]
\centering
\caption{実験の条件}
\label{tab:e-cond}
\footnotesize
\begin{tabularx}{\linewidth}{@{}l>{\raggedright\arraybackslash}X@{}}
\toprule
固定 & 既定のまま使ったものは「既定の設定」と書く \\
\midrule
比較 & 何を \\
 & どの範囲で \\
 & 計何 run \\
\midrule
掃引 & 何を振るか \\
 & どの範囲で \\
 & 計何 run \\
\bottomrule
\end{tabularx}
\end{table}

\subsection{取った特徴量}\label{sec:e-feat}

成績のほかに記録した量を，表にする（1種類1行）。取る理由（仮定）はREPのmdに書く。

\begin{table}[H]
\centering
\caption{取った特徴量}
\label{tab:e-feat}
\footnotesize
\begin{tabularx}{\linewidth}{@{}l>{\raggedright\arraybackslash}X@{}}
\toprule
種類 & 記録した量 \\
\midrule
種類1 & 量の名前を並べる \\
種類2 & 量の名前を並べる \\
\bottomrule
\end{tabularx}
\end{table}

% ============================================================
\section{結果}\label{sec:e-result}
% ============================================================

実験で観測した事実だけを，簡潔に書く。段落の頭に，内容を表す短い見出しを太字で置く
（\textbf{動作}，\textbf{規模感}，\textbf{手法の比較}など）。原因の説明や評価の言葉はここに書かない。
最後に，実験の前に決めた達成条件（REPのmd）を満たしたかを，1文で書く。

% 探索型は，結果で終わる。結果から生まれた問いは，PDFに書かず，REPのmdに記録する（2026-10-10 の決定）。

\end{document}
